\documentclass{article}

\usepackage{amsmath,amssymb}
\usepackage{xurl}
\usepackage[hidelinks]{hyperref}

% Shared commands for notes in docs/paper-gaps/.

\DeclareUnicodeCharacter{2080}{\ifmmode{}_{0}\else\textsubscript{0}\fi}
\DeclareUnicodeCharacter{2081}{\ifmmode{}_{1}\else\textsubscript{1}\fi}
\DeclareUnicodeCharacter{2082}{\ifmmode{}_{2}\else\textsubscript{2}\fi}
\DeclareUnicodeCharacter{2099}{\ifmmode{}_{n}\else\textsubscript{n}\fi}

\newcommand{\Lean}{\textsc{Lean}}
\ExplSyntaxOn
\NewDocumentCommand{\leanid}{m}{
  \ifmmode
    \textsf{\detokenize{#1}}
  \else
    \group_begin:
    \urlstyle{sf}
    \tl_set:Nn \l_tmpa_tl {#1}
    \tl_replace_all:Nnn \l_tmpa_tl {\_} {_}
    \exp_args:NV \path \l_tmpa_tl
    \group_end:
  \fi
}
\ExplSyntaxOff
\providecommand{\tr}{\operatorname{tr}}
\newcommand{\tnleanIssueBase}{https://github.com/LionSR/TNLean/issues/}
\newcommand{\tnleanPrBase}{https://github.com/LionSR/TNLean/pull/}
\newcommand{\tnleanDocBase}{https://sirui-lu.com/TNLean/docs/}
\newcommand{\ghissue}[1]{\href{\tnleanIssueBase#1}{GitHub issue \##1}}
\newcommand{\ghpr}[1]{\href{\tnleanPrBase#1}{PR \##1}}
\newcommand{\doclink}[1]{\href{\tnleanDocBase#1.html}{\path{#1}}}

% Dirac notation and the identity operator, as in blueprint/src/macros/common.tex.
% \providecommand keeps note-local definitions authoritative.
\usepackage{dsfont}
\providecommand{\ket}[1]{|#1\rangle}
\providecommand{\bra}[1]{\langle #1|}
\providecommand{\braket}[2]{\langle #1 | #2 \rangle}
\providecommand{\ketbra}[2]{|#1\rangle\!\langle #2|}
\providecommand{\Id}{\mathds{1}}
\providecommand{\one}{\mathds{1}}
\providecommand{\C}{\mathbb{C}}
\providecommand{\MN}[1]{M_{#1}(\C)}
\providecommand{\MD}{M_D(\C)}

% External referencing for paper sources.  The build helper writes sanitized
% auxiliary files under build/paper-aux/ and excludes duplicated source labels.
\usepackage{xr}

% Import a generated paper auxiliary file with a caller-supplied label prefix.
% The candidate paths support builds from docs/paper-gaps/, the repository root,
% and build/paper-aux/ itself.
\newcommand{\importpaperaux}[2]{%
  \IfFileExists{../../build/paper-aux/#2.aux}{%
    \externaldocument[#1]{../../build/paper-aux/#2}%
  }{%
    \IfFileExists{build/paper-aux/#2.aux}{%
      \externaldocument[#1]{build/paper-aux/#2}%
    }{%
      \IfFileExists{#2.aux}{%
        \externaldocument[#1]{#2}%
      }{}%
    }%
  }%
}

% General external reference macros
\newcommand{\paperref}[2]{\ref{#1-#2}}
\newcommand{\papereqref}[2]{(\ref{#1-#2})}

% CPSV16 source (1606.00608 / MPDO-22-12-17-2.tex) external document and shortcuts
\importpaperaux{cpsv16-}{1606.00608/MPDO-22-12-17-2-tnlean}
\newcommand{\cpsvref}[1]{\paperref{cpsv16}{#1}}
\newcommand{\cpsveqref}[1]{\papereqref{cpsv16}{#1}}

% Verdict marker: \gapnote{<kind>}{<status>}. Kind is the policy
% classification (clarification, local-correction, scope-restriction,
% unfaithful, false-source, open-gap); status is open, wip (elimination
% underway), resolved, or historical. Machine-read by the site index;
% prints a verdict line.
\newcommand{\gapnote}[2]{\par\noindent\textbf{Verdict:} #1 (#2).\par}


\title{Physical-Isometry Transport for the Blocked MPDO Fixed Point}
\author{}
\date{2026-07-12}

\begin{document}
\maketitle
\gapnote{scope-restriction}{resolved}

\section{Source construction and present scope}

Proposition~\texttt{3to5} of Ref.~\cite{Cirac2016MPDO_arXiv} constructs
channels $\mathcal T$ and $\mathcal S$ from a supplied factorization of the
physical sector; the relevant argument appears in Appendix~C.2,
lines~1510--1563 and 1821--1825. At lines~1821--1825, the construction
composes each channel with an overlapping copy to obtain the two-site-blocked
channels $\mathcal T^2$ and $\mathcal S^2$. These formulas apply to one
factorization of the entire input tensor.

Condition~(iv) of Theorem~4.9 of Ref.~\cite{Cirac2016MPDO_arXiv} instead
provides a factorization for each representative in the basis of normal
tensors (BNT) and imposes no constraint on its unnormalized repeated-copy
weights. The scalar presentation with weights $1$ and $1/2$ satisfies
condition~(iv), but its two-site block fails Definition~4.1 of
Ref.~\cite{Cirac2016MPDO_arXiv}. Consequently, the single-factorization
transport theorem is a conditional auxiliary result. Combining outer physical
supports cannot extend it to the literal implication
(iv)$\Rightarrow$(v), which is refuted.

The outer-support construction still applies to the viable source proposition
\texttt{prop2to5}, which begins with condition~(ii) of
Ref.~\cite{Cirac2016MPDO_arXiv}. In the preceding Case-II derivation, zero
correlation length (ZCL) makes the repeated-copy weights independent of the
copy label, after which their common value is absorbed into the normalization.
Only after this normalization can the channel pairs from
Proposition~\texttt{3to5} be combined over mutually orthogonal outer physical
supports.

Let $\mathcal K$ denote the original tensor, let $\widehat{\mathcal K}$ denote
the same tensor in the sector coordinates selected by the physical isometry,
and let $\operatorname{blockTwo}$ denote two-site blocking. The
sector-coordinate theorem proves the fixed-point identities for
\begin{align}
    \operatorname{blockTwo}(\widehat{\mathcal K}).
    \label{eq:cpgsv17_mpdo_blocked_rfp_physical_transport_sector_block}
\end{align}
\footnote{Formal declaration:
\leanid{NeighboringTraceFactorization.blockTwo_sectorCoordinateTensor_isRFPViaTS}.}
Its conclusion is exact in the selected sector coordinates, but it does not
transport the two channels back to
$\operatorname{blockTwo}(\mathcal K)$. It is therefore a scope-restricted
intermediate result for one supplied factorization of the entire tensor.

\section{Transport theorem}

For $n\in\{2,4\}$ and an admissible auxiliary input $X$, let
$\mathcal K_n(X)$ and $\widehat{\mathcal K}_n(X)$ denote the corresponding
$n$-site physical operators in the original and sector coordinates,
respectively. In the formal sector factorization, the physical coordinate map
is a square matrix $U$ satisfying
\begin{align}
    U^\dagger U
        &= I.
    \label{eq:cpgsv17_mpdo_blocked_rfp_physical_transport_isometry}
\end{align}
Finite dimensionality and squareness then give
\begin{align}
    UU^\dagger
        &= I.
    \label{eq:cpgsv17_mpdo_blocked_rfp_physical_transport_coisometry}
\end{align}
Thus $U$ is unitary.

With the tensor indices arranged according to the right-associated convention
used in the channel construction, the transport theorem establishes
\begin{align}
    \widehat{\mathcal K}_n(X)
        &= U^{\otimes n}\mathcal K_n(X)(U^\dagger)^{\otimes n},
        \qquad n\in\{2,4\}.
    \label{eq:cpgsv17_mpdo_blocked_rfp_physical_transport_unitary_closure}
\end{align}
Let $\widehat{\mathcal T}^2$ and $\widehat{\mathcal S}^2$ denote the blocked
channels in the sector coordinates. Transporting these channels through the
unitary congruences in
Eq.~\ref{eq:cpgsv17_mpdo_blocked_rfp_physical_transport_unitary_closure}
preserves complete positivity and trace preservation. The two
sector-coordinate closure identities therefore yield the required identities
for the original blocked tensor.

\section{Discharge}

The transport theorem formalizes the two-site and four-site unitary closure
relations, transports both channels to the original physical matrix spaces,
and proves
\begin{align}
    \operatorname{IsRFPViaTS}
        (\operatorname{blockTwo}(\mathcal K)).
    \label{eq:cpgsv17_mpdo_blocked_rfp_physical_transport_original_block_rfp}
\end{align}
\footnote{Formal declaration:
\leanid{NeighboringTraceFactorization.blockTwo_isRFPViaTS}.}
This theorem discharges the physical-coordinate transport gap for one supplied
factorization of the entire tensor. It does not prove the raw implication
(iv)$\Rightarrow$(v) in Theorem~4.9 of
Ref.~\cite{Cirac2016MPDO_arXiv}, which is false.

For the source implication (ii)$\Rightarrow$(v), after absorption of the
common weight derived from ZCL, the theorem supplies the channel pair for each
factorization. The projector-controlled assembly over all BNT representatives
is formalized in \ghissue{6632}. The remaining source obligation is the
factorization for each representative, documented in
\path{docs/paper-gaps/cpgsv17_pf_rank_one.tex} and tracked in
\ghissue{6775}.

\bibliographystyle{alpha}
\bibliography{references}

\end{document}
